

\clearpage
\section{Problema de investigación}
¿Qué es?
Es una descripción clara, objetiva y contextualizada de la realidad que genera una necesidad, conflicto o desafío.
¿Cómo redactarlo?
Describe hechos concretos (datos, cifras, observaciones).
Ubica el problema en un contexto: ¿Dónde ocurre? ¿A quién afecta?
Evita generalizaciones o juicios de valor.

\subsection{Situación Problemática}
En este acápite se expondrá de manera clara, integral, articulada y sin ambigüedades la situación problemática (características del problema), haciendo énfasis en la situación de las variables contenidas en el tema y su relación. En ese marco se identificará el vacío de conocimiento, las dudas, las contradicciones que se requiere resolver y las nuevas perspectivas a estudiar.
Finalmente, se expresará el enunciado del problema, como una pregunta clara y concreta, que exprese las variables principales y su relación. La problemática debe incluir al menos tres aspectos: (a) antecedentes científicos fundamentados con referencias bibliográficas actualizadas de la problemática y del problema específico a resolver; (b) exponer de manera clara, precisa y concreta el vacío del conocimiento de la problemática indicada; (c) incluir la pregunta central a la cual la investigación propuesta responderá, esta pregunta servirá de base para construir el objetivo general del proyecto; además de la pregunta central, se incluirán las preguntas específicas de las cuales se derivarán los objetivos específicos e hipótesis. El enunciado del problema se puede formular con preguntas del siguiente tipo: ¿qué efecto?, ¿en qué condiciones...?, ¿cuál es la probabilidad de...? .
Para concretar la problemática es necesario un dominio del marco referencial\footnote{\href{https://unl.edu.ec/investigacion/guias-investigacion}{Guía de TIC}}; es decir, se debe tener conocimientos teóricos de la temática a investigar. La información que fundamenta la problemática debe estar respaldada con citas bibliográficas. Para estructurar la problemática se puede hacer uso de técnicas que permitan el análisis y organización de la información colectada, analizada y sistematizada, entre las principales se anotan:

\subsection{Pregunta de investigación}
Se pueden utilizar cualquiera de ellas, dependiendo del tema que se esté tratando, el éxito está en seleccionar la adecuada.
Diagrama causa-efecto
Lluvia de ideas
Matriz problemática
Matriz de necesidades
Árbol de problemas
\subsection{Justificación}
En esta sección se resalta la importancia de la investigación, considerando las siguientes premisas: exponer el por qué es necesario realizar la investigación, indicar el aporte del proyecto al avance de conocimientos, cómo apoyará a la solución de problemas, demostrar para qué y quiénes se investigará (p.ej., demandas de grupos sociales específicos, prioridades institucionales) . Además, se hará constar algunos aspectos como: la relación con los Objetivos de Desarrollo Sostenible (ODS), objetivos de desarrollo nacional o local, políticas estatales o institucionales, la relación con las líneas de investigación a nivel general y de las carreras o programas a la que pertenece el estudiante\footnote{\href{https://unl.edu.ec/investigacion}{Dirección de Investigación}}.